```tex
\chapter*{Resumen}
\addcontentsline{toc}{chapter}{Resumen}
Las crecientes exigencias de trazabilidad y eficiencia en la manufactura agroindustrial contrastan con la realidad de numerosos laboratorios universitarios y pequeñas empresas, donde maquinaria electromecánica heredada opera como cajas negras sin conectividad digital. En el Laboratorio de Procesos de la Universidad Indoamérica, una máquina dosificadora de granos controlada por un PLC de arquitectura cerrada y circuito de potencia de 220V carecía de cualquier capacidad de monitorización o interacción externa, limitando su valor académico y operativo. Ante esta problemática, el presente trabajo se propuso desarrollar un gemelo digital interactivo y predictivo de Nivel 3 sin modificar el cableado interno del sistema original. Para ello se diseñó e implementó un bypass ciberfísico no invasivo basado en microcontroladores ESP32, que captura pasivamente las señales de sensores infrarrojos, ultrasónico y de presión mediante optoacopladores, y las transmite vía MQTT hacia un SCADA en Node-RED y un simulador cinemático tridimensional en Webots. La lógica de control se modeló como una Máquina de Estados Finitos programada en Python, asegurando la réplica fiel del ciclo de dosificación sin interferir con el autómata original. La metodología siguió un modelo en V adaptado a sistemas ciberfísicos, incorporando prácticas de desarrollo ágil y pruebas con inyección de fallos. Los resultados cuantitativos obtenidos durante más de 76 minutos de operación validan la efectividad del sistema: la Efectividad General de los Equipos (OEE) promedio alcanzó un 91.85\% con una desviación estándar de ±0.46\%, la latencia de telemetría MQTT se situó en 0.61 ms y, ante un fallo de presión inyectado, el tiempo medio de recuperación fue de 52 segundos, con una disponibilidad del 92.92\% durante la contingencia. Todo el sistema se materializó con un presupuesto de 243.20 USD empleando exclusivamente herramientas de código abierto. Se concluye que la arquitectura de bypass ciberfísico constituye una solución viable, de bajo costo y alta replicabilidad para dotar de inteligencia predictiva a maquinaria legacy, convirtiendo un activo opaco en un recurso didáctico y de investigación alineado con los paradigmas de la Industria 4.0.

\textbf{Palabras clave:} Gemelo Digital, Sistema Ciberfísico, Bypass Ciberfísico, Máquina de Estados Finitos, Industria 4.0

\chapter*{Abstract}
\addcontentsline{toc}{chapter}{Abstract}
The growing demands for traceability and efficiency in agribusiness manufacturing contrast starkly with the reality of many university laboratories and small enterprises, where legacy electromechanical machinery operates as black boxes without digital connectivity. In the Process Laboratory of Universidad Indoamérica, a grain dosing machine controlled by a closed-architecture PLC and a 220V power circuit lacked any monitoring or external interaction capability, limiting its academic and operational value. In response to this challenge, the present work aimed to develop a Level 3 interactive and predictive digital twin without modifying the system's internal wiring. A non-invasive cyber-physical bypass was designed and implemented using ESP32 microcontrollers, which passively capture signals from infrared, ultrasonic and pressure sensors through optocouplers and transmit them via MQTT to a Node-RED SCADA and a three-dimensional kinematic simulator in Webots. The control logic was modeled as a Finite State Machine programmed in Python, ensuring faithful replication of the dosing cycle without interfering with the original controller. The methodology followed a V-Model adapted to cyber-physical systems, incorporating agile development practices and fault injection testing. Quantitative results obtained over more than 76 minutes of operation validate the system's effectiveness: the average Overall Equipment Effectiveness (OEE) reached 91.85\% with a standard deviation of ±0.46\%, MQTT telemetry latency was 0.61 ms, and during an injected pressure fault the mean time to recovery was 52 seconds, with an availability of 92.92\% during the contingency. The entire system was realized with a budget of 243.20 USD using exclusively open-source tools. It is concluded that the cyber-physical bypass architecture constitutes a viable, low-cost, and highly replicable solution to endow legacy machinery with predictive intelligence, transforming an opaque asset into a didactic and research resource aligned with Industry 4.0 paradigms.

\textbf{Keywords:} Digital Twin, Cyber-Physical System, Cyber-Physical Bypass, Finite State Machine, Industry 4.0

\chapter*{Resumen Ejecutivo}
\addcontentsline{toc}{chapter}{Resumen Ejecutivo}

La transformación digital del sector agroindustrial requiere que incluso los activos más antiguos puedan integrarse en ecosistemas ciberfísicos. No obstante, en laboratorios universitarios y pequeñas y medianas empresas de Ecuador, máquinas electromecánicas heredadas, como la dosificadora de granos del Laboratorio de Procesos de la Universidad Indoamérica, operan bajo controladores propietarios sin puertos de comunicación modernos. Esta condición las convierte en cajas negras que impiden el registro de variables de proceso, la detección temprana de anomalías y el desarrollo de competencias en Industria 4.0. El presente proyecto aborda esta brecha proponiendo un gemelo digital de Nivel 3, predictivo e interactivo, que se acopla externamente a la máquina sin alterar su lógica original.

El objetivo general fue desarrollar dicho gemelo digital mediante una integración ciberfísica no invasiva basada en un bypass con microcontroladores ESP32, el motor de simulación Webots, una Máquina de Estados Finitos programada en Python y una capa de comunicaciones MQTT con supervisión en Node-RED. El sistema debía emular fielmente la cinemática del ciclo industrial y permitir la validación de métricas de rendimiento y resiliencia frente a fallos. Para alcanzarlo se definieron tres objetivos específicos: el levantamiento diagnóstico de la lógica secuencial del activo físico, el diseño e implementación del modelo ciberfísico completo, y la evaluación mediante la inyección controlada de perturbaciones.

Metodológicamente se adoptó el modelo en V adaptado a sistemas ciberfísicos, complementado con técnicas ágiles propias del desarrollo de software. La fase inicial consistió en una ingeniería inversa de la máquina original mediante trazados de continuidad y análisis de señales eléctricas, respetando en todo momento la política de no manipulación interna. A partir de este diagnóstico se diseñó un panel de control externo en gabinete IP54 que alberga dos ESP32, sensores infrarrojos, un sensor ultrasónico, un transductor de presión, un módulo de relés de ocho canales y un servomotor. El bypass se instrumentó con optoacopladores PC817 para garantizar aislamiento galvánico entre el circuito de potencia de 220V y la electrónica de bajo voltaje, capturando pasivamente las señales de los actuadores sin interferir en el comando original del PLC. La arquitectura de comunicación se basó en el protocolo MQTT, con un broker Mosquitto desplegado en un entorno Linux; los tópicos se segmentaron en \texttt{maquina/sensor/\#} y \texttt{maquina/actuador/\#}, transmitiendo payloads en formato JSON a una frecuencia de muestreo de 100 ms. El SCADA se construyó en Node-RED, ofreciendo un dashboard con indicadores en tiempo real de OEE, estados de la máquina y registro histórico en base de datos SQLite. El modelo cinemático tridimensional se desarrolló en FreeCAD y se importó a Webots, donde los actuadores se controlan mediante nodos SliderJoint y HingeJoint gobernados por una FSM implementada en Python que replica el ciclo de alimentación, dosificación y sellado. La comunicación bidireccional entre el mundo virtual y el físico cierra así el lazo ciberfísico sin vulnerar la integridad del controlador legado.

Los resultados cuantitativos respaldan la eficacia del sistema. Durante cinco pruebas con lotes de 60 a 200 fundas y un tiempo acumulado de operación de 76.11 minutos, se obtuvo una Efectividad General de los Equipos promedio de 91.85\% con una desviación estándar de apenas ±0.46\%, lo que evidencia alta estabilidad. La disponibilidad fue del 100\% en condiciones normales y del 92.92\% durante la prueba con fallo inyectado. La latencia media de la telemetría MQTT fue de 0.61 ms, cumpliendo holgadamente el requisito de diseño de menos de 100 ms. La capacidad predictiva se comprobó al inyectar una caída de presión neumática: el protocolo de recuperación se ejecutó en 52 segundos, métrica de MTTR que antes del proyecto era imposible de cuantificar. Todo el prototipo se construyó con un presupuesto ejecutado de 243.20 USD, utilizando exclusivamente herramientas de código abierto.

Es necesario reconocer ciertas limitaciones. La métrica OEE no incorpora el factor de calidad por no disponerse de un sistema de pesaje integrado, y las pruebas se realizaron en un ambiente controlado de laboratorio con un único tipo de fallo simulado. Pese a ello, el impacto del proyecto es significativo: se ha transformado un activo opaco en un sistema monitorizado con trazabilidad total, se ha generado una metodología de bypass ciberfísico replicable en otras máquinas legacy y se ha creado un caso de estudio de bajo costo que puede transferirse directamente a entornos académicos y a pequeñas empresas agroindustriales. Como recomendaciones ejecutivas, se plantea la incorporación de una celda de carga para completar el cálculo integral de OEE, la extensión de la arquitectura hacia gemelos de Nivel 4 con inteligencia artificial prescriptiva, y la formalización de alianzas con el sector productivo para validar el sistema en condiciones industriales reales. El proyecto demuestra que la Industria 4.0 no exige necesariamente renovar el parque de maquinaria, sino que puede comenzar con intervenciones inteligentes, no intrusivas y de bajo costo.

\chapter*{Executive Summary}
\addcontentsline{toc}{chapter}{Executive Summary}

The digital transformation of the agribusiness sector requires that even the oldest assets be integrated into cyber-physical ecosystems. However, in university laboratories and small and medium-sized enterprises in Ecuador, legacy electromechanical machines, such as the grain dosing machine at the Process Laboratory of Universidad Indoamérica, operate under proprietary controllers without modern communication ports. This condition turns them into black boxes that prevent process variable logging, early anomaly detection, and the development of Industry 4.0 skills. The present project addresses this gap by proposing a Level 3, predictive and interactive digital twin that couples externally to the machine without altering its original logic.

The overall objective was to develop this digital twin through a non-invasive cyber-physical integration based on a bypass with ESP32 microcontrollers, the Webots simulation engine, a Finite State machine programmed in Python, and an MQTT communication layer with Node-RED supervision. The system had to faithfully emulate the kinematics of the industrial cycle and enable the validation of performance and resilience metrics against failures. To achieve this, three specific objectives were defined: diagnostic mapping of the physical asset's sequential logic, complete design and implementation of the cyber-physical model, and evaluation through controlled disturbance injection.

Methodologically, the V-Model adapted to cyber-physical systems was adopted, complemented by agile software development techniques. The initial phase consisted of reverse engineering the original machine using continuity traces and electrical signal analysis, always respecting the internal no-manipulation policy. Based on this diagnosis, an external control panel was designed in an IP54 enclosure housing two ESP32s, infrared sensors, an ultrasonic sensor, a pressure transducer, an eight-channel relay module, and a servomotor. The bypass was instrumented with PC817 optocouplers to ensure galvanic isolation between the 220V power circuit and the low-voltage electronics, passively capturing actuator signals without interfering with the original PLC command. The communication architecture was based on the MQTT protocol, with a Mosquitto broker deployed on a Linux environment; topics were segmented into \texttt{maquina/sensor/\#} and \texttt{maquina/actuator/\#}, transmitting JSON payloads at a sampling rate of 100 ms. The SCADA was built in Node-RED, providing a dashboard with real-time OEE indicators, machine states, and historical logging in an SQLite database. The three-dimensional kinematic model was developed in FreeCAD and imported into Webots, where actuators are controlled by SliderJoint and HingeJoint nodes governed by an FSM implemented in Python that replicates the feeding, dosing, and sealing cycle. The bidirectional communication between the virtual and physical worlds thus closes the cyber-physical loop without compromising the integrity of the legacy controller.

Quantitative results support the system's effectiveness. During five tests with batches of 60 to 200 bags and an accumulated operation time of 76.11 minutes, an average Overall Equipment Effectiveness of 91.85\% was obtained with a standard deviation of only ±0.46\%, indicating high stability. Availability was 100\% under normal conditions and 92.92\% during the fault-injected test. The average MQTT telemetry latency was 0.61 ms, comfortably meeting the design requirement of less than 100 ms. Predictive capability was proven when a pneumatic pressure drop was injected: the recovery protocol executed in 52 seconds, an MTTR metric that was impossible to quantify before the project. The entire prototype was built with an executed budget of 243.20 USD, using exclusively open-source tools.

Certain limitations must be recognized. The OEE metric does not incorporate the quality factor due to the absence of an integrated weighing system, and the tests were conducted in a controlled laboratory environment with a single type of simulated fault. Despite this, the project's impact is significant: an opaque asset has been transformed into a monitored system with full traceability, a replicable cyber-physical bypass methodology has been generated for other legacy machines, and a low-cost case study has been created that can be directly transferred to academic settings and small agribusiness companies. As executive recommendations, the incorporation of a load cell to complete the integral OEE calculation is proposed, as well as extending the architecture towards Level 4 digital twins with prescriptive artificial intelligence, and formalizing partnerships with the productive sector to validate the system under real industrial conditions. The project demonstrates that Industry 4.0 does not necessarily require renewing the machine park, but can start with intelligent, non-intrusive, and low-cost interventions.
```